\section{Invertibility Modulo \texorpdfstring{$m$}{m}}

\begin{proposition}[Compatibility]
Let $a_1, a_2, b_1, b_2 \in \mathbb{Z}$. If $a_1 \equiv a_2 \pmod{m}$ and $b_1 \equiv b_2 \pmod{m}$, then:
\begin{enumerate}
    \item $a_1 + b_1 \equiv a_2 + b_2 \pmod{m}$
    \item $a_1b_1 \equiv a_2b_2 \pmod{m}$
\end{enumerate}
\end{proposition}
\begin{proof}
    \leavevmode
\begin{itemize}
    \item \textbf{Addition:} Given $a_2 \equiv a_1 \pmod{m}$ and $b_2 \equiv b_1 \pmod{m}$, $\implies \exists k_1, k_2 \in \mathbb{Z} : a_2 - a_1 = k_1 m \text{ and } b_2 - b_1 = k_2 m$.
    Adding the equations yields $(a_2 - a_1) + (b_2 - b_1) = (k_1 + k_2) m$.
    Rearranging the terms: $(a_2 + b_2) - (a_1 + b_1) = (k_1 + k_2) m$.
    Since $k_1 + k_2 \in \mathbb{Z}$, this shows $m \mid ((a_2 + b_2) - (a_1 + b_1))$, so $a_1 + b_1 \equiv a_2 + b_2 \pmod{m}$.
    \item \textbf{Multiplication:} We express $a_2 = a_1 + k_1 m$ and $b_2 = b_1 + k_2 m$.
    Multiplying them gives:
    $$ a_2b_2 = (a_1 + k_1 m)(b_1 + k_2 m) = a_1b_1 + a_1 k_2 m + b_1 k_1 m + k_1 k_2 m^2 $$
    Factoring out $m$: $a_2b_2 - a_1b_1 = m(a_1 k_2 + b_1 k_1 + k_1 k_2 m)$.
    Let $K = a_1 k_2 + b_1 k_1 + k_1 k_2 m$. Since $K \in \mathbb{Z}$, $m \mid (a_2b_2 - a_1b_1)$, so $a_1b_1 \equiv a_2b_2 \pmod{m}$. \qedhere
\end{itemize}
\end{proof}

\begin{center}
    \vspace{10pt}
    \line(1,0){450}
    \vspace{25pt}
\end{center}

\begin{proposition}
The set of integers modulo $m$ under addition, $(\mathbb{Z}/m\mathbb{Z}, +)$, forms a \textbf{commutative group}.
\end{proposition}

\begin{center}
    \vspace{10pt}
    \line(1,0){450}
    \vspace{25pt}
\end{center}

\begin{proposition}
When considering both addition and multiplication, $(\mathbb{Z}/m\mathbb{Z}, +, \cdot)$ forms a \textbf{commutative ring}.
\end{proposition}

\begin{center}
    \vspace{10pt}
    \line(1,0){450}
    \vspace{25pt}
\end{center}

\begin{theorem}
If an inverse exists, it is unique.
\end{theorem}
\begin{proof}
Assume an element $[a]$ has two inverses, $[a']$ and $[a'']$. By definition, $[a] \cdot [a'] = [1]$ and $[a] \cdot [a''] = [1]$. Since multiplication is associative and $[1]$ is the identity element, we can write:
$$ [a'] = [a'] \cdot [1] = [a'] \cdot ([a] \cdot [a'']) $$
$$ [a'] = ([a'] \cdot [a]) \cdot [a''] = [1] \cdot [a''] = [a''] $$
Therefore, $[a'] = [a'']$, $\implies$ the inverse is unique.
\end{proof}

\begin{center}
    \vspace{10pt}
    \line(1,0){450}
    \vspace{25pt}
\end{center}

\begin{theorem}[Characterization of Invertibility]
An element $[a] \in \mathbb{Z}/m\mathbb{Z}$ is invertible $\iff \gcd(a, m) = 1$.
\end{theorem}
\begin{proof}
An element $[a]$ is invertible $\iff \exists u \in \mathbb{Z} : a u \equiv 1 \pmod{m}$.
This congruence is equivalent to the linear Diophantine equation $a u - 1 = -m v$ for some $v \in \mathbb{Z}$, which can be rewritten as $a u + m v = 1$.
By Bézout's Identity, this equation has an integer solution $\iff \gcd(a, m) \mid 1$.
Since $\gcd(a, m) \in \mathbb{N}$, this condition is satisfied $\iff \gcd(a, m) = 1$. \qedhere
\end{proof}

\begin{center}
    \vspace{10pt}
    \line(1,0){450}
    \vspace{25pt}
\end{center}

\begin{proposition}
The set of units under multiplication, $(\mathbb{Z}/m\mathbb{Z})^*$, forms a \textbf{commutative group} (the group of units modulo $m$).
\end{proposition}

\begin{center}
    \vspace{10pt}
    \line(1,0){450}
    \vspace{25pt}
\end{center}

\begin{proposition}
$\forall p$ prime number, the commutative ring $(\mathbb{Z}/p\mathbb{Z}, +, \cdot)$ is a \textbf{field}, commonly denoted $\mathbb{F}_p$.
\end{proposition}

\begin{center}
    \vspace{10pt}
    \line(1,0){450}
    \vspace{25pt}
\end{center}

\begin{corollary}
Let $p \in \mathbb{N}$ be a prime so that $p > 2$. The equation $x^2 \equiv 1 \pmod{p}$ in $\mathbb{F}_p$ has exactly two distinct solutions: $x \equiv 1 \pmod{p}$ and $x \equiv p-1 \pmod{p}$ (i.e., $x \equiv -1 \pmod{p}$).
\end{corollary}
\begin{proof}
The equation $x^2 \equiv 1 \pmod{p}$ is equivalent to $p \mid (x^2 - 1)$, or $p \mid (x - 1)(x + 1)$.
Since $p$ is prime, by the Prime Divisor Property:
\begin{itemize}
    \item Either $p \mid (x - 1)$, which implies $x \equiv 1 \pmod{p}$.
    \item Or $p \mid (x + 1)$, which implies $x \equiv -1 \equiv p-1 \pmod{p}$.
\end{itemize}
These two solutions are distinct because $1 \not\equiv -1 \pmod{p}$ for $p > 2$. \qedhere
\end{proof}

\begin{center}
    \vspace{10pt}
    \line(1,0){450}
    \vspace{25pt}
\end{center}